% =====================================================================
% Section 7 --- Conclusion (merged draft; supersedes (7)discussion.tex
% and the stub (8)conclusion.tex)
% ---------------------------------------------------------------------
% Job: return to the thesis; position against BF as the frictionless pole;
%      carry the closure lineage (Sen, Taylor, Rose, FIDELIO); state the
%      limitations; close with practical guidance. NO bridge-discovery
%      claims -- the bridge is measurement. Kuhnian frame kept minimal.
% Mindmap: N03, N02D.
% Status: first full draft (2026-10-22). Absorbs the two paragraphs the
%      dropped Discussion would have carried (positioning, limitations).
% NUMBERS TO INSERT once the Sec. 5 prose is final: sign-reversal count and
%      the multiplier range across the five labor closures (Tables 4/5).
% RETIRED (2026-10-24): "the Delta Pi/Y range from the financing wedge (Sec. 5.2)".
%      No financing-wedge section was written -- Sec. 5 carries only aggregate and
%      sectoral results -- and the number does not exist: the kernel prices at
%      marginal cost, so Delta Pi/Y is identically 0.00 in every plain financing
%      cell and is tabulated only once a pricing door is switched on (Tables 7
%      and E). Do not re-open without a new run generation.
% =====================================================================
\section{Conclusion}\label{sec:conclusion}

Input--output models of a socio-ecological transformation report multipliers, employment
effects and fiscal bills as if the short run were a matter of arithmetic:
invert a matrix, read off the change. What is key for estimating contemporary relationships (e.g. when using environmentally extended input-output tables) is, maybe, the least interesting step when focusing on an exogenous demand shock (with a transformative purpose). The latter requires assumptions on wage closures and financing regimes which requires considering competing theories of what the short run is, and which mechanisms underpin our monetary system. Moreover, it can strongly impact the magnitude of the results, and occasionally also their sign.

Our contribution is not a new model but an illustrative comparison focusing on a practical purpose. One core framework substantiates fifteen scenarios whose differences are entirely attributable to the assumptions that define them. Read this way,
the matrix does two things at once. It identifies where the closure choice
is redundant and shows where it matters most, either in terms of magnitude (with the no-reallocation model and Leontief closure as poles) or in terms of creating structurally distinct economies (e.g., sectoral models as prone to producing relative price changes).

This result is, hence, strongly embedded within the existing literature.
The accounting with intermediate inputs of \citet{Hulten.1978} and its
modern descendants \citep{Baqaee.2021,Acemoglu.2012} describe the
frictionless pole of our matrix --- the limit in which sectoral labor supply
is infinitely elastic and reallocation carries no cost. Our five closures represent  canonical answers of a long debate --- \citet{Sen1963} on the demarcation of neoclassical and neo-Keynesian closures, \citet{Taylor1990} on structuralist closures, \citet{Rose.1995} on their translation into computable form ---
carried into an accounting structure that transformation research already
uses. FIDELIO \citep{Rocchi.2025} operationalizes one of these closures; the
matrix makes that closure an explicit modeling decision rather than a
background assumption.

Three limitations bound the claims. First, labor is the only production factor in the model, so our results speak to the labor-market channel of the short run and are somewhat silent on the capital- and investment-driven side of adjustment. Nor is the factor itself differentiated: a skill differential --- skilled and unskilled
labor as separate factors --- would allow the same closure question to be asked about
the composition of labor demand rather than only about its level, which is where much of
the recent evidence on the transition is located (Section~\ref{sec:io-bridge}). Our
single-factor design cannot answer that question, and it is the extension we would
pursue first. Second, the welfare metric is dispersion-blind --- real consumption deflated by $\mathcal{P}$ says nothing about who within the household sector bears
the adjustment. Third, the horizon is static: we compare short-run equilibria, not adjustment paths. In this context, an interesting vantage point for a dynamic reinterpretation of our efforts is given by \citet{McGregor.1996}, who shows how a capacity-constrained economy, where investment responds to demand, will converge into a model with demand-determined quantities.

%TODO: hier eventuell nochmal die toolbox für das applied modeling einstreuen – siehe auch remark in section 1? JK: habs mal gemaht; müsste man vllt noch schärfen oder verlinken.

For the applied modeler and transformation researcher, our paper provides a first, illustrative menu of closure choices and a corresponding open toolbox to put them into practice. For any interested reader, it might be refreshing to see how the older lesson survives the newer formalism. The short run is a modeling choice, and the choice carries the paradigm with it. Which short run --- and whose money --- remains the question; the numbers answer it only after the assumptions have been named.
